%% kapa: exam=1 round=21 year=2010 period=1 form=A subject=economics subject_ko=경제원론 seq=12 number=52
%% source: Q-Net 감정평가사 기출문제 (cst00309) artlSeq=5024181 file=extracted/1008/5024181/2010년 제21회 감정평가사 1차시험 기출문제/1교시.hwp
%% conversion: kordoc 4.24.51 -> tools/gen1.py (auto); review=proofread-sonnet; flags=none
\begin{question}{52}
X재와 Y재를 소비하는 甲의 소비량에 따른 한계효용이 다음 표와 같다. X재의 가격이 10, Y재의 가격이 20일 때 효용극대화 조건이 충족되는 甲의 소비묶음은?

{\small\centering
\begin{tabular}{|l|c|c|c|c|c|c|}
\hline
소비량 & 1 & 2 & 3 & 4 & 5 & 6 \\
\hline
X재의 한계효용 & 10 & 9 & 8 & 7 & 6 & 5 \\
\hline
Y재의 한계효용 & 10 & 8 & 6 & 5 & 3 & 2 \\
\hline
\end{tabular}\par}

\begin{choices}
  \item X = 1, Y = 4
  \item X = 2, Y = 3
  \item X = 3, Y = 2
  \item X = 5, Y = 3
  \item X = 6, Y = 1
\end{choices}
\end{question}
\answer{5}
